\documentclass[11pt]{article}

\usepackage[margin=1.05in]{geometry}
\usepackage{amsmath,amssymb}
\usepackage{graphicx}
\usepackage{booktabs}
\usepackage{microtype}
\usepackage{xcolor}
\usepackage[hidelinks]{hyperref}
\usepackage{caption}
\usepackage{subcaption}
\usepackage{natbib}
\usepackage{url}

\captionsetup{font=small,labelfont=bf}
\bibliographystyle{plainnat}

\newcommand{\dd}{\mathcal{D}}
\newcommand{\code}[1]{\texttt{\small #1}}

\title{\bf What Makes a Mission Hard?\\[2pt]
\large Defining, Measuring, and Learning Difficulty in a\\
Game-Scale Gravitational Navigation Environment}

\author{Kevin Horecka\thanks{Correspondence: \texttt{kevin.horecka@gmail.com}.
Code, levels and every measurement reported here:
\url{https://github.com/kevroy314/spaceship-navigator}.}}

\date{\today}

\begin{document}
\maketitle

\begin{abstract}
We present \emph{spacenav}, a simulated navigation environment in which a
thrusting spacecraft must visit a sequence of moving targets inside a
gravitational field produced by many bodies. The environment is cheap to
simulate and easy to state, and yet we find that the ordinary question
``how hard is this mission?'' resists a single answer in three distinct ways.
\emph{First}, difficulty is not one quantity: a mission can be hard because a
simple mental model of the gravitational field is wrong, or because the window
for acting is narrow, and these two failure modes require opposite fixes.
We give a closed-form, simulation-free score, $\dd = \eta / A$, that captures
the first and can be \emph{inverted} to construct missions at a chosen
difficulty, and we argue that it must never be collapsed with the second.
\emph{Second}, difficulty is hard to measure even once defined. Our own
published conclusion that a class of two-stop missions was infeasible turned
out to be an artifact of twelve correlated samples; at thirty-two samples the
same class is flown sixteen percent of the time. We report this and two other
self-corrections as findings about the measurement problem rather than
footnotes.
\emph{Third}, difficulty interacts with what the learner is told. Across twenty
training runs at five random seeds per condition, we show that adding four
hand-designed physical readouts to the agent's observation \emph{costs} it
roughly ten percentage points of success, and that the entire cost is
attributable to the two readouts that are algebraically redundant given the
rest of the observation, while the two that carry genuinely new physics are
free. We release the environment, the levels, the trained policies and the
analysis code.
\end{abstract}

\section{Introduction}

A spacecraft with limited fuel must reach a target that is itself orbiting.
The physics is Newtonian and deterministic; the simulation runs in milliseconds;
a competent undergraduate can write the equations of motion from memory. And yet
if you ask ``is this particular mission hard?'' the question turns out to have no
single well-posed answer, and the ways in which it fails to have one are, we
think, more interesting than the environment itself.

This paper is about that failure, using one concrete environment as the subject.
\emph{Spacenav} is a two-dimensional, game-scale gravitational navigation task:
a ship with a steerable engine and a finite fuel tank must fly through or match
velocity with a short sequence of targets, some of which are attached to moving
bodies, inside a field generated by up to thirty-two masses. Episodes last 150
simulated seconds at a 15\,Hz control rate, and the action space is the nine
combinations a human could press on a gamepad. The environment is designed as a
testbed for reinforcement learning, but the difficulties we report are not
specific to learning: a search-based trajectory optimiser meets most of them too.

We make three claims.

\paragraph{Difficulty needs at least two axes, and one of them has a closed form.}
The natural mesoscale model of a gravitational field --- ``only the body I am
orbiting matters'', known in astrodynamics as \emph{patched conics} --- is wrong
by a measurable amount at every point in space. We define $\eta$ as the fraction
of the local gravitational acceleration that such a one-body model discards, and
$A$ as the ship's fuel budget expressed in units of the impulse gravity will
deliver over the remaining flight. Their ratio $\dd = \eta / A$ is the share of
the tank that flying the wrong model would cost. It requires no simulation, so
millions of candidate missions can be scored for free, and it inverts: given a
target difficulty one can solve for the fuel that places a mission there. But
$\dd$ says nothing about \emph{timing}, and a mission can be trivial to plan and
impossible to execute. We keep the two separate throughout, and show what
collapsing them would hide (Section~\ref{sec:difficulty}).

\paragraph{Measuring difficulty is itself hard, and we got it wrong in public.}
The honest way to ask whether a mission class is feasible is to let a strong
optimiser attempt it. We did that, with three search budgets, and reported that
a class of two-stop missions never succeeded. The result was wrong: the three
budgets had been run against the \emph{same} twelve missions, so they were not
three independent attempts, and the true success rate is about sixteen percent,
at which drawing zero out of twelve has probability $0.13$. We describe this and
two further self-corrections in Section~\ref{sec:measuring}, because a benchmark
whose difficulty estimates are this fragile is a fact about the benchmark.

\paragraph{What you tell the agent changes what it can learn, in a direction that
is easy to get backwards.}
We instrumented the environment with four physically meaningful readouts ---
$\eta$, the fuel-versus-gravity ratio, their quotient $\dd$, and a local tidal
stretching rate --- on the hypothesis that an agent told about the structure of
its own difficulty would learn faster. The opposite happened. Across five
training seeds per condition, the four channels cost $3.1$ to $10.1$ percentage
points of final success depending on the evaluation tier, all four differences
significant. Splitting the panel in half shows the cost is \emph{entirely}
carried by the two channels that are algebraic combinations of quantities the
agent already observes; the two that carry information not otherwise present are
statistically indistinguishable from supplying nothing at all
(Section~\ref{sec:instruments}).

Throughout, we have tried to report what we measured rather than what we
expected, including where the two differed. The environment's own development
history is littered with elegant, confident and wrong conclusions, each caught
only by a measurement, and we have found that record more useful than the
successes.

\section{Related work}
\label{sec:related}

\paragraph{Astrodynamics and the coarse model.}
The idea that one body dominates a spacecraft's local dynamics is the
\emph{patched-conic approximation}, standard in mission design and accurate
precisely when the quantity we call $\eta$ is small. Its domain of validity is
usually described by the Hill radius or the Laplace sphere of influence. We
instead use the \emph{equal-pull surface}, the radius $d = a\sqrt{m/M}$ at which
a secondary of mass $m$ at semi-major axis $a$ pulls as hard as its primary of
mass $M$; this is roughly half a Hill radius and is where $\eta$ peaks.
Learning-based guidance for low-thrust transfers has been studied by
\citet{izzo2021} among others, who report that policy networks trained on
optimal transfers recover near-optimal control with small optimality gaps; our
setting differs in being game-scale, partially observed and adversarially
generated rather than drawn from a mission catalogue.

\paragraph{Chaos, and whether it is present.}
Our initial design intent was for the environment to be chaotic, since chaotic
regimes are where coarse mental models fail most interestingly. The standard
route to constructed chaos in a planetary context is the Chirikov resonance
overlap criterion \citep{chirikov1979} and Wisdom's $2/7$ law for the chaotic
zone around a secondary \citep{wisdom1980}, refined by \citet{deck2013}. We
rejected this route on timescale grounds and verified the rejection empirically
(Section~\ref{sec:chaos}): resonance overlap governs evolution over \emph{many}
orbital periods, whereas an episode here is a fraction of one. For basin
structure we use the uncertainty exponent $\alpha$ of
\citet{grebogi1983}, under which the fraction of uncertain initial conditions
scales as $f(\varepsilon)\sim\varepsilon^{\alpha}$, with $\alpha\!\approx\!1$
indicating a smooth, resolvable boundary and $\alpha\!\approx\!0$ a fractal one;
basin entropy \citep{daza2016} is a related diagnostic.

\paragraph{Why chaos would matter for a learner.}
\citet{wang2023} show that a reinforcement learning objective is
$\alpha$-H\"older with $\alpha = \log(1/\gamma)/\lambda$ for a system with
maximal Lyapunov exponent $\lambda$ and discount $\gamma$, and becomes fractal
--- losing any useful descent direction --- when $\lambda$ exceeds
$\log(1/\gamma)/\Delta t$. At our discount of $\gamma = 0.999$ and control
interval $\Delta t = 1/15$\,s this threshold is
$\lambda_{\mathrm{crit}} = 0.0150\,\mathrm{s}^{-1}$. We measured against it
rather than assuming either side of it.

\paragraph{Policy compression and execution precision.}
That ``how simple is the rule'' and ``how precisely must it be timed'' are
different quantities is argued from two directions:
policy compression in human decision making \citep{lai2024} and policy
information capacity as a task-difficulty measure \citep{furuta2021}. A
bang-bang controller on a razor-thin switching surface is a one-bit policy on
an impossible trigger. Our execution-margin thresholds are taken from the human
motor literature: approximately $20$\,ms onset jitter for a practised discrete
action \citep{repp2005} and a closed-loop correction ceiling near $3$\,Hz
\citep{franklin2019}.

\paragraph{Evaluation practice.}
Our statistical protocol follows \citet{agarwal2021} and
\citet{colas2018}: separate training-seed variance from episode variance, use
enough seeds to say anything about the former, and prefer interval estimates.
This discipline is the reason several results in this paper are reported as
null. For context on the scale of modern environment-step budgets we cite
\citet{kinetix2024} and \citet{simbav2}, both of which make clear that our
training budget is small by current standards --- a limitation we return to in
Section~\ref{sec:limits}.

\section{The environment}
\label{sec:env}

\subsection{Physics}

The simulation is two-dimensional with gravitational constant $G = 1$, so each
``mass'' is really a gravitational parameter $GM$. Lengths are game units (u) and
times are game seconds (s); a typical level spans 1000--3000\,u. Bodies attract
one another and are integrated with a kick-drift-kick leapfrog scheme at a
substep of $1/60$\,s. Within a body's radius the field falls off linearly, as
for a uniform sphere, which keeps the acceleration finite everywhere and makes
self-interaction exactly zero.

The ship is a \emph{test particle}: it feels every body but moves none. This is
the single most consequential design decision in the environment. Because the
ship exerts no force, a level's body trajectories are completely determined by
its initial snapshot, which means (i) the entire future of the field can be
precomputed and replayed exactly, (ii) a browser client can reproduce the
physics from a $\sim$1.5\,KB snapshot by integrating the bodies itself, and
(iii) a level is reproducible from a seed with no stored trajectory data.

The ship's acceleration is the sum of gravity, atmospheric drag where a body has
an atmosphere, and thrust. Thrust is held constant across a control tick. Four
physics substeps make one control tick, so the agent and a human player both act
at $15$\,Hz, or every $66.7$\,ms.

\subsection{Missions, hazards and termination}

A \emph{mission} is a short tour: between one and six targets, each either a
point in free space, a station, or a point attached to a moving body. Each
target carries an arrival radius (default $12$\,u) and a velocity tolerance;
a flythrough target accepts any speed, while a \emph{rendezvous} target requires
matching to within $3$\,u/s, which is a materially harder manoeuvre. A tour may
be ordered or order-free. The ship carries a finite tank, expressed as seconds
of full-throttle burn, and a per-mission engine acceleration; both are sampled
per mission from ranges chosen so that a ship can never simply ignore gravity
and fly straight at the target.

Hazards are radiation from luminous bodies, atmospheric heating, debris fields
and radiation belts. An episode ends on arrival at the final target, collision
with a solid body, destruction by accumulated damage, fuel exhaustion with no
remaining way to arrive, departure past $1.35$ times the level radius, or
timeout at $150$\,s ($2250$ control ticks).

\subsection{Observation and action}

The observation is ego-centric and set-structured, so that the same weights run
on a level with any number of bodies. It comprises a fixed ship-and-tour block
(velocity, local gravity vector, relative target position and velocity, fuel
fraction, remaining $\Delta v$ against the tour's estimated requirement, health,
time remaining, and the objective weights), one token per body (relative
position and velocity, distance, the body's contribution to local acceleration,
mass, radius, atmosphere, luminosity, and a one-hot kind), and one token per
hazard zone. Bodies and zones are embedded token-wise and pooled two ways:
multi-head cross-attention with the ship's own features as the query, and a
masked max-pool.

The action space is deliberately the set a person can press: turn
$\in\{\text{left},\text{none},\text{right}\}$ crossed with throttle
$\in\{\text{off},\text{gentle},\text{full}\}$, nine discrete actions in total.
Turning is arcade-style with no angular inertia.

\subsection{Level families and pools}

Levels are generated procedurally in twelve named families, from
\code{sol\_like} (a dominant star with well-separated planets) through
\code{binary\_star}, \code{trinary}, \code{star\_cluster}, \code{bh\_cluster},
\code{debris\_maze} and \code{rogue\_flyby}. A generated level is accepted only
if it remains well-behaved --- no collisions, no ejections, bounded energy drift
--- for $900$\,s, six times a single episode. Accepted levels are stored as
snapshots at $20$\,s intervals so that an episode may begin at any of several
epochs in the level's evolution.

Four pools are used: a training pool of $18{,}000$ levels, a validation pool of
$2{,}700$ levels drawn from the same families, a holdout pool of $1{,}200$
levels from families never seen in training, and a long-horizon pool of $1{,}800$
levels carrying sixteen snapshots for experiments at extended episode clocks.
Pools are reproducible from seeds.

\section{Defining difficulty}
\label{sec:difficulty}

\subsection{The coarse model and what it discards}

The mental model a person brings to a gravitational field is almost always the
one-body model: identify the dominant attractor, treat its field as the whole
field, and reason about conic sections. We instantiate exactly this model in
code. The function \code{gravity\_at(..., topk=k)} restricts the ship's felt
gravity to the $k$ strongest attractors at its current position and discards the
rest; $k=1$ is patched conics, and $k=0$ is the truth. Bodies always move
truthfully --- only the test particle's felt force is restricted --- so a plan
made inside the coarse model can be flown in the true world and the discrepancy
observed directly.

This gives the first axis of difficulty a natural definition. Let
$a_{\mathrm{full}}$ be the true local acceleration and $a_{\mathrm{dom}}$ the
acceleration from the single strongest body. Define
\begin{equation}
\eta \;=\; \frac{\lVert a_{\mathrm{full}} - a_{\mathrm{dom}} \rVert}
{\lVert a_{\mathrm{full}} \rVert},
\end{equation}
the fraction of the local field that a one-body picture throws away. Near zero,
patched conics is right; near one, it is lying.

\subsection{Authority, and the ratio that matters}

$\eta$ alone is not difficulty, because a ship with a large enough engine can
simply absorb the error. Let $\Delta v_{\mathrm{tank}}$ be the mission's fuel
expressed as a velocity budget, $\lVert g \rVert$ the local gravitational
acceleration and $T$ the remaining flight time. Define the \emph{authority}
\begin{equation}
A \;=\; \frac{\Delta v_{\mathrm{tank}}}{\lVert g\rVert\, T},
\end{equation}
the tank measured in units of the impulse gravity itself will deliver over the
flight. $A > 1$ means the ship can out-push the field; $A < 1$ means it cannot.

The quantity we use to score a mission is the ratio
\begin{equation}
\boxed{\;\dd \;=\; \frac{\eta}{A}\;}
\end{equation}
which has a direct reading: \emph{the share of the fuel tank that flying the
wrong model would cost}. Both ingredients are available in closed form from the
level snapshot and the mission specification, so $\dd$ requires no simulation at
all; thousands of candidate missions can be scored essentially for free, and
expensive planning is spent only on missions already known to sit where they are
wanted.

$\dd$ also \emph{inverts}. Solving for the tank,
\begin{equation}
\Delta v_{\mathrm{tank}} \;=\; \frac{\eta \,\lVert g \rVert\, T}{\dd},
\end{equation}
so a mission can be placed on a chosen difficulty rung by solving rather than by
generate-and-test. One practical caution: when retuning a mission onto a rung we
only ever \emph{add} fuel. Reaching a high rung by starving the tank produces an
infeasible mission dressed up as a demanding one, which is a different thing.

\paragraph{Calibration.} An ordinary mission in this environment sits at
$\dd \approx 0.0014$, with $\eta = 0.0036$ and $A = 3.8$. Gravity supplies about
a quarter of the tank's worth of impulse over a flight, but a \emph{single} body
explains $99.6\%$ of the field. That is why patched conics works here, and why
breaking it requires deliberate construction rather than luck.

\paragraph{Validation.} We compared a planner restricted to the one-body model
against a planner with the true field, on missions pooled by $\dd$ band, 24
missions per band. Below $\dd = 0.1$ the two tie ($32\%$ versus $42\%$ arrival,
Fisher exact $p = 0.74$); above it they separate sharply ($13\%$ versus $54\%$,
$p = 0.005$), with the coarse planner failing every mission on the top rung. We
note explicitly that the \emph{per-rung} figures, at $n = 8$, are not
significant; only the pooled comparison is. This is the first of several places
where our own early reporting outran its sample size.

\subsection{Where to put a target: the equal-pull surface}

To raise $\eta$ constructively one needs a place to stand. The literature's
obvious candidate is the chaotic zone around a secondary, delimited by the
resonance-overlap criterion. We tried it and rejected it. Resonance overlap
describes what happens over many orbital periods; an episode here is a fraction
of one, so asymptotic chaos never gets time to act. The measurement was blunt:
placing targets inside the Wisdom band lifted $\eta$ from $0.0026$ to $0.0172$,
a sixfold increase that is still negligible, because a few Hill radii from a
secondary is an $\eta \approx 0.02$ desert.

What governs an episode this short is instantaneous force competition. The
relevant surface is where a secondary's pull equals its primary's, at
$d = a\sqrt{m/M}$ --- about half a Hill radius.
Figure~\ref{fig:eta}(a) shows $\eta$ measured on a shell around the
second-most-massive body of $297$ levels spanning nine families, taken as a
median over $24$ directions at each radius. It rises from $0.09$ at
$0.3\times$ the equal-pull radius to a peak of $0.43$ at $1.0\times$, then falls
to $0.24$, $0.066$ and $0.020$ at $2\times$, $4\times$ and $8\times$. Target
placement in the contested-mission generator is therefore confined to the
$0.8$--$1.3\times$ shell.

We note that the peak value depends strongly on how one samples it. Along the
line joining primary and secondary the two pulls oppose, the denominator nearly
cancels and $\eta$ exceeds $1$; the angular median reported here is the
conservative statistic and is the one the generator is tuned against.

The larger lever is family choice (Figure~\ref{fig:eta}(b)). Median $\eta$ over
the whole playable area is $0.0005$ for \code{debris\_maze} and $0.0027$ for
\code{sol\_like}, against $0.35$ for \code{binary\_star} and $0.53$ for
\code{star\_cluster} --- a spread of three orders of magnitude. Systems of
comparable masses are contested everywhere; hierarchical systems are contested
only in thin shells that must be aimed at deliberately.

\begin{figure}[t]
\centering
\includegraphics[width=\textwidth]{figures/eta.pdf}
\caption{(a) The fraction $\eta$ of local gravity that a one-body model
discards, against distance from a secondary in units of the equal-pull radius
$a\sqrt{m/M}$. Median over $24$ directions on each shell, across $297$ levels
from nine families; the band is the interquartile range. The coarse model is
worst almost exactly where the two pulls balance. (b) Median $\eta$ over the
whole playable area, by level family. Families in orange are contested
everywhere; the rest are contested only inside a thin shell.}
\label{fig:eta}
\end{figure}

\subsection{The second axis, and why it must stay separate}
\label{sec:twoaxis}

$\dd$ measures whether a coarse \emph{model} suffices. It says nothing about
whether a correct plan can be \emph{executed}. We therefore band every mission on
two axes and never reduce them to one number:

\begin{itemize}
\setlength\itemsep{1pt}
\item \textbf{Model scale} --- the coarsest $k$ in \code{gravity\_at(topk=k)}
under which a planner's trajectory still arrives when flown in the true world.
\item \textbf{Execution margin} --- timing slack against measured human limits:
$\sim\!20$\,ms motor onset jitter, $\sim\!200$\,ms visual reaction time, and a
$3$\,Hz closed-loop correction ceiling. Our control tick is $66.7$\,ms, already
coarser than a trained hand, which makes the comparison generous rather than
rigged.
\end{itemize}

\begin{figure}[t]
\centering
\includegraphics[width=0.60\textwidth]{figures/bands.pdf}
\caption{The two difficulty axes kept apart. A mission's position on the
horizontal axis says whether a coarse one-body model suffices; its position on
the vertical axis says whether the timing window is within human reach.
\emph{Calculated} and \emph{precision} are different kinds of impossible and
need opposite repairs, which is why a single scalar difficulty would merge
exactly the two cases a designer most needs to tell apart.}
\label{fig:bands}
\end{figure}

The resulting bands (Figure~\ref{fig:bands}) are \emph{trivial}, \emph{human}, \emph{calculated}
(intuition fails but computation succeeds), \emph{precision} (no hand can hold
it) and \emph{opaque} (nothing we have succeeds; rejected as unfair). The reason
to keep the axes apart is that \emph{calculated} and \emph{precision} are
different kinds of impossible and need opposite repairs: widen a
\emph{precision} mission's timing window and it becomes \emph{calculated};
simplify a \emph{calculated} mission's field and it becomes \emph{human}.
A single scalar difficulty would merge exactly the two cases a level designer
most needs to tell apart.

\section{Is the environment chaotic?}
\label{sec:chaos}

Several of our design threads --- fractal basin boundaries, a ``lottery''
difficulty band, and a worry that a fractal objective would leave a policy
gradient with no descent direction --- presupposed that the field was chaotic
on episode timescales. It is not, and establishing that changed the plan.

The concern with sharpest teeth was the H\"older result of \citet{wang2023}: at
$\gamma = 0.999$ and $15$\,Hz, a maximal Lyapunov exponent above
$\lambda_{\mathrm{crit}} = 0.0150\,\mathrm{s}^{-1}$ would make the objective
fractal. Our measured predictability horizon --- the time at which one control
tick's worth of thrust error grows past the $12$\,u arrival tolerance --- is
$52$--$74$\,s across eight families, uncomfortably close to the $67$\,s
corresponding to that threshold.

The resolution is that the divergence is \emph{ballistic, not exponential}. The
discriminator is the decay of $d\ln(\mathrm{sep})/dt$, which falls as $1/t$ for
linear growth and stays flat for exponential growth; pure linear growth predicts
a $150$\,s to $75$\,s slope ratio of exactly $0.50$, and we measure $0.33$--$0.64$
with median $\approx 0.56$. Fitting
$\mathrm{sep} = a\,t\,e^{\lambda t}$ to divide out the unavoidable ballistic
factor gives $\lambda = 0.0045$--$0.0086\,\mathrm{s}^{-1}$ (90th percentile
$\approx 0.010$), roughly a factor of two below $\lambda_{\mathrm{crit}}$; only
the tail of \code{star\_cluster} exceeds it. A Lyapunov time of $\sim\!140$\,s
against $150$\,s episodes is about one e-folding per episode, or $\sim\!2000$
control steps per e-folding, whereas the successful chaotic-control results in
the literature sit at $10$--$100$ steps per e-folding. We are comfortably inside
the benign regime.

The uncertainty exponent told the same story the hard way. Our first estimates
of $\alpha$ were computed from too few samples: with 96 perturbation samples the
boundary-flip counts were small integers (0, 1, 3, 5, 5), one fitted exponent of
$1.297$ was a two-point slope between two and six flips, and bootstrapped
confidence intervals straddled the threshold separating fractal from resolvable.
Zero of five exponents was trustworthy, and we had reported a fractal/resolvable
split that the data did not support. After raising the sample count eightfold,
$\alpha$ measures $0.53$--$1.02$: resolvable, consistent with the Lyapunov
result. We also added an explicit \code{measurable} flag so that an
unmeasurable exponent propagates as \emph{unknown} rather than as \emph{fractal}.

Three loose ends resolve at once. The environment is \emph{quasi-integrable} on
episode timescales. Two difficulty bands that were supposed to capture chaotic
missions contain zero instances, because there is no chaos to capture. And a
previously published claim of ours, that a training plateau was caused by a
fractal objective, is retracted: that plateau needs a different explanation.

\section{Measuring difficulty empirically}
\label{sec:measuring}

\subsection{The planner as yardstick, and the solvability gate}

Closed-form $\dd$ scores a mission's model demand, but it cannot tell you whether
a mission is flyable at all. For that we use a cross-entropy-method trajectory
optimiser as a reference pilot: piecewise-constant control segments over the
episode, seeded from a scripted reflex controller and from an explicit
zero-thrust candidate, iteratively refined toward an arrival objective. A
mission counts as \emph{solvable} if the optimiser arrives, or if any member of
an enumerated one- or two-burn family arrives. Missions that nothing can fly are
excluded from difficulty labelling: a level nobody can fly earns no label, since
its measured difficulty would describe the probe rather than the level.

Including the zero-thrust candidate matters more than it sounds. Throttle is
parameterised through a sigmoid, so an all-zero-throttle plan is essentially
never drawn by the search; on missions whose optimal solution is to coast and let
gravity do the work, the optimiser was therefore failing on the easiest levels
in the environment until the candidate was added explicitly.

\subsection{A retraction, and what it says about the benchmark}
\label{sec:retraction}

We previously reported that two-stop missions --- tours with two sequential
targets, typically with at least one rendezvous --- were \emph{infeasible} at the
$150$\,s episode clock rather than merely hard. The evidence was that twelve
such missions were attempted under three search budgets, including one with
sixteen times the trajectory evaluations, and none succeeded, with timeout as
the dominant failure.

That conclusion was wrong, and the way it was wrong is instructive. The three
budgets were run against the \emph{same} twelve missions, so they were close to
perfectly correlated and did not constitute three independent chances to observe
a success. Re-running with thirty-two missions per condition and a modestly
larger search gives a success rate of $5/32 \approx 16\%$ at the same clock. At
that rate, observing zero successes in twelve attempts has probability
$(1-5/32)^{12} = 0.13$ --- an unremarkable outcome, not evidence of
impossibility.

Extending the episode clock while holding control resolution fixed at $3.1$\,s
per segment gives the picture in Figure~\ref{fig:clock}. Arrivals rise from
$5/32$ at $150$\,s to $13/32$ at $600$\,s and then stop; timeouts fall
monotonically from $18$ to $5$; and the episodes that a longer clock frees up
become hazard endings about as often as they become arrivals: crashes,
destructions and fuel exhaustions go from $9$ at the shortest clock to $16$ at
the longest, though not monotonically ($9$, $12$, $10$, $16$).
More time buys more exposure. The useful ceiling for this mission class is
around $600$\,s, and the binding constraint beyond it is hazard exposure rather
than the clock.

\begin{figure}[t]
\centering
\includegraphics[width=0.46\textwidth]{figures/clock.pdf}
\caption{Thirty-two two-stop missions flown by the reference optimiser at four
episode clocks, with control resolution held at $3.1$\,s per segment so that a
longer clock buys more segments rather than coarser ones. Timeouts fall
monotonically; arrivals peak near $600$\,s; hazard endings rise throughout. An
earlier report of ours, based on twelve correlated samples at the $150$\,s
clock, concluded that this mission class was infeasible.}
\label{fig:clock}
\end{figure}

\subsection{What the episode clock is, structurally}

The episode clock is a global constant compared against a tick counter, not a
per-mission field. This is a limitation we discovered while running the sweep
above --- each clock required a separate process with the constant patched at
import --- and it is the reason the environment cannot presently run long
missions for play while measuring difficulty on a short, fixed clock. We note it
because it is exactly the kind of structural detail that silently constrains
which experiments are cheap.

\section{Learning}
\label{sec:learning}

\subsection{Agent and training}

We train with Proximal Policy Optimization \citep{schulman2017} in a fully
vectorised setup where rollout, advantage estimation and all minibatch epochs
compile into a single accelerated call. The policy is the actor-critic described
in Section~\ref{sec:env}: embedding width $128$, four attention heads,
$285{,}322$ parameters in total.

Hyperparameters: $1024$ parallel environments $\times\ 128$ steps $=131{,}072$
transitions per update, four epochs, eight minibatches, learning rate
$3\times10^{-4}$ annealed linearly to a tenth of that, $\gamma = 0.999$,
generalised-advantage $\lambda = 0.95$, clip $0.2$, entropy coefficient annealed
from $0.01$ to $0.002$, value coefficient $0.5$, gradient-norm clip $0.5$. Runs
are $3000$ updates, or $393$ million environment steps.

A curriculum ramps three things over training: the probability that a leg is a
rendezvous rather than a flythrough ($0.1 \to 0.45$), the maximum tour length
($1 \to 3$), and the proportion of missions whose targets are placed in the
contested equal-pull shell ($0 \to 0.35$, held at zero until the ship can fly at
all). The fuel budget and engine acceleration start generous and tighten to
their full ranges.

\subsection{Evaluation suites}

Four frozen suites of approximately $600$ missions each are held fixed across all
runs:
\emph{easy} (one flythrough target, roomy tank, strong engine, three hierarchical
families), \emph{medium} (one or two targets, some rendezvous, tighter budget,
all training families), \emph{hard} (two or three targets, half of them
rendezvous, a budget usually insufficient for a direct flight so that gravity
must do some of the work), and \emph{blind} (families never seen in training,
with order-free tours $30\%$ of the time). Six hundred missions per suite gives
roughly two percentage points of binomial noise on a success rate.

Each suite also carries a reference score from a scripted reflex pilot ---
proportional navigation with gravity compensation --- evaluated on the identical
missions.

\subsection{Results}

Table~\ref{tab:main} gives final performance for the best configuration against
the reflex pilot. Success is the mean of each run's last three checkpoints,
averaged over five training seeds.

\begin{table}[t]
\centering
\small
\begin{tabular}{lrrrr}
\toprule
& Easy & Medium & Hard & Blind \\
\midrule
Learned policy (5 seeds) & $89.4 \pm 1.5$ & $41.7 \pm 3.0$ & $10.7 \pm 1.2$ & $30.2 \pm 2.3$ \\
Scripted reflex pilot    & $85.3$ & $34.7$ & $7.6$ & $31.8$ \\
\midrule
Difference               & $+4.1$ & $+7.0$ & $+3.1$ & $\mathbf{-1.6}$ \\
\bottomrule
\end{tabular}
\caption{Percentage of missions arrived. The learned policy beats the scripted
pilot on three suites and loses on \emph{blind}, the only suite drawn from
families never seen in training and the one where tour order is sometimes the
pilot's to choose.}
\label{tab:main}
\end{table}

Two observations deserve emphasis. First, the learned policy \emph{loses} to a
hand-written reflex controller on the held-out families. Second, the reflex
pilot is a weaker yardstick than it appears: on a comfortable-tank mission set
it arrives $63.3\%$ of the time, and its dominant failure is \emph{collision} at
$21.6\%$, not fuel exhaustion ($4.7\%$) or timeout ($7.8\%$). Any margin quoted
against this opponent is correspondingly soft, a caveat we attach to every such
comparison. Absolute success rates are unaffected.

\section{What the agent should be told}
\label{sec:instruments}

\subsection{Motivation and design}

Having defined $\eta$, $A$ and $\dd$ as the quantities that describe a mission's
difficulty, the natural question is whether an agent told those quantities
learns better. We added four channels to the observation: $\eta$; a
log-compressed authority $A$; a log-compressed demand $\dd$; and a local tidal
stretching rate, the square root of the tidal gradient, which measures how fast
nearby trajectories are pulled apart. These are intended as a cockpit: the
quantities a well-instrumented pilot would read.

The channels can be zeroed individually or together, leaving the observation
\emph{shape} unchanged, so that an ablated run remains comparable with a full
one and no checkpoint is invalidated. We verified separately that toggling them
does not perturb the dynamics.

\subsection{The four channels cost roughly ten points}

Figure~\ref{fig:ablation}(a) and Table~\ref{tab:instr} give the result at five
training seeds per arm. Supplying all four channels costs $3.1$ to $10.1$
percentage points of final success, with Welch's $t$-test on seed means
significant on all four suites and effect sizes between $1.8$ and $3.0$. We
emphasise that these are tests on \emph{seed means}: the quoted uncertainties are
training-seed standard deviations, not episode-level binomial intervals, which
are much smaller and would overstate confidence.

\begin{table}[t]
\centering
\small
\begin{tabular}{lrrrr}
\toprule
& Easy & Medium & Hard & Blind \\
\midrule
No channels      & $89.4 \pm 1.50$ & $41.7 \pm 3.04$ & $10.7 \pm 1.24$ & $30.2 \pm 2.27$ \\
$\eta$, stretch  & $90.4 \pm 1.81$ & $41.9 \pm 2.30$ & $11.1 \pm 1.03$ & $29.2 \pm 1.23$ \\
budget, demand   & $86.6 \pm 1.22$ & $33.2 \pm 2.61$ & $6.2 \pm 1.68$  & $22.7 \pm 2.52$ \\
All four         & $86.3 \pm 1.87$ & $31.6 \pm 5.25$ & $5.8 \pm 3.39$  & $22.2 \pm 2.99$ \\
\midrule
$p$ (none vs.\ $\eta$,stretch)   & $0.40$ & $0.90$ & $0.60$ & $0.42$ \\
$p$ (all vs.\ budget,demand)     & $0.77$ & $0.59$ & $0.81$ & $0.79$ \\
$p$ ($\eta$,stretch vs.\ budget,demand) & $\mathbf{0.0064}$ & $\mathbf{0.0005}$ & $\mathbf{0.0011}$ & $\mathbf{0.0022}$ \\
\bottomrule
\end{tabular}
\caption{Percentage of missions arrived, mean $\pm$ standard deviation over five
training seeds per arm, $3000$ updates each. Welch's $t$-test on seed means. The
two informative channels are statistically indistinguishable from supplying
nothing; the two redundant channels are indistinguishable from supplying all
four.}
\label{tab:instr}
\end{table}

\begin{figure}[t]
\centering
\includegraphics[width=\textwidth]{figures/ablation.pdf}
\caption{(a) Final success for every training seed, with no instrument channels
against all four. Dots are individual seeds, bars are one standard deviation,
the heavy tick is the mean. (b) The four-way split. In every suite the
no-channel and $\eta$,stretch arms coincide, and the budget,demand and all-four
arms coincide below them.}
\label{fig:ablation}
\end{figure}

\subsection{The cost is carried entirely by the redundant channels}

The explanation was visible in the observation code before it was measured. The
observation already carries the local gravity vector (hence
$\lVert g \rVert$), the product of fuel and engine acceleration (hence
$\Delta v_{\mathrm{tank}}$), and the fraction of the episode remaining (hence
$T$). Therefore $A = \Delta v_{\mathrm{tank}} / (\lVert g \rVert T)$ is an
algebraic combination of three quantities the agent already has, and partly
duplicates an existing ratio of remaining to required $\Delta v$; and
$\dd = \eta/A$ is redundant given $\eta$ and $A$. By contrast $\eta$ requires the
dominant-body decomposition, which is not otherwise computable from the
observation, and the tidal stretching rate is not present at all.

This generated a prediction: splitting the panel along the redundancy line should
put the entire cost on the redundant half. It does. Supplying only $\eta$ and
the stretching rate is statistically indistinguishable from supplying nothing
($p = 0.40$ to $0.90$ across suites); supplying only $A$ and $\dd$ is
indistinguishable from supplying all four ($p = 0.59$ to $0.81$); and the
contrast between the two halves is significant on every suite with effect sizes
from $2.4$ to $3.6$.

We draw a narrow conclusion and resist a broad one. Narrowly: in this
environment, at this network size and step budget, appending algebraically
redundant features to an observation measurably degrades a policy, while
appending features carrying genuinely new physics is free. We do \emph{not} have
evidence that the informative channels help --- $\eta$ and stretch are nominally
above the no-channel arm on three suites and below on one, none significantly.
And we have not isolated the mechanism. Four extra input dimensions against
$285{,}322$ parameters is a capacity story; a ratio with a division in it is
heavy-tailed even after log compression, which is a conditioning story. Both
remain untested.

\section{What a single number hides}
\label{sec:ghosts}

Evaluation reports one success rate per checkpoint, computed from a greedy
policy that takes the most probable action. That number conceals both the
policy's own stochasticity and the enormous variation between individual
missions. Figure~\ref{fig:ghosts} makes both visible: every checkpoint of one
training run flying one fixed mission sixteen times, with actions \emph{sampled}
rather than maximised, for $31 \times 16 = 496$ trajectories per mission.

\begin{figure}[p]
\centering
\includegraphics[width=\textwidth]{figures/ghosts.pdf}
\caption{Four missions from the evaluation suites, each flown $496$ times by the
same training run: every saved checkpoint, sixteen sampled rollouts each,
coloured by the training update at which the attempt was flown. Grey discs are
bodies with their orbits traced faintly; green rings are targets, drawn with
their own paths where they are attached to moving bodies. Early policies (dark)
spray across the field; late policies (bright) converge into tight bundles.
\emph{Medium, task 1} converges onto a single arrival corridor.
\emph{Medium, task 0} converges onto a corridor that mostly misses.
\emph{Hard, task 0} converges onto a tight shuttle between its first two targets
and never reaches the third, which sits at upper left: an agent confidently
executing a plan that cannot complete the mission.}
\label{fig:ghosts}
\end{figure}

Two things follow. First, sampled success is far below greedy success --- the
\emph{easy} mission shown arrives $52\%$ of the time when sampled against $89\%$
for the greedy policy on that suite --- so the headline number describes a
behaviour the policy exhibits only when its stochasticity is removed. Second,
the spread \emph{between missions within a suite} is larger than any effect we
measure between training configurations: $60\%$, $52\%$, $8\%$ and $0\%$ on the
four missions shown. The instrument-channel effect of Section~\ref{sec:instruments}
is seven to ten points; two missions drawn from the same suite differ by fifty.
Any claim about this environment that rests on a handful of missions is
therefore fragile by construction, which is the lesson of
Section~\ref{sec:retraction} restated in pictures.

Figure~\ref{fig:curves} shows the corresponding learning curves.

\begin{figure}[t]
\centering
\includegraphics[width=\textwidth]{figures/curves.pdf}
\caption{Evaluation success against training update for one seed of each arm.
Mid-run checkpoints are noisy enough that reading an ablation before it finishes
is actively misleading; we did so twice during this work and drew the wrong
conclusion both times.}
\label{fig:curves}
\end{figure}

\section{Threats to validity}
\label{sec:threats}

\paragraph{The reference pilot is weak.} Every margin quoted against the
scripted pilot is softened by its $21.6\%$ collision rate
(Section~\ref{sec:learning}). Absolute rates are unaffected.

\paragraph{The hard suite is partly unflyable.} The reference optimiser arrives
on only about $16\%$ of two-stop missions at the measurement clock, and the
\emph{hard} suite draws from that family. An agent scoring $10.7\%$ against a
$3000$-trajectory-per-decision optimiser scoring $16\%$ is a different statement
from an agent scoring $10.7\%$ against a solvable task.

\paragraph{Training budget.} At $393$ million environment steps our runs are
small relative to contemporary work at the $10^{11}$ scale \citep{kinetix2024},
and at $285{,}322$ parameters the network is small relative to current
general-purpose agents \citep{simbav2}. Results here should be read as
statements about this operating point, not asymptotically.

\paragraph{Seeds.} Five training seeds per arm meets the minimum we set
ourselves, following \citet{agarwal2021}, but it is a minimum. Interval
estimates on effect sizes at $n=5$ remain wide.

\paragraph{Two-dimensional physics.} The environment is planar. Several
phenomena that matter in real mission design --- plane changes, out-of-plane
perturbations --- cannot occur.

\section{Reproduction}
\label{sec:repro}

All code, level generators, evaluation suites and analysis scripts are at
\url{https://github.com/kevroy314/spaceship-navigator}. The environment is
implemented in JAX; training runs reported here were executed on a single
consumer GPU. The repository additionally contains a numbered, append-only series
of decision records documenting the reasoning behind each design choice,
including the ones that were later retracted; most of the self-corrections
reported in this paper are traceable to a specific record and the measurement
that forced it.

Representative commands:

\begin{quote}\small
\begin{verbatim}
python scripts/build_pools.py                  # regenerate level pools from seeds
python scripts/train.py --run myrun --updates 3000 --seed 0
SPACENAV_INSTRUMENTS=eta,stretch python scripts/train.py --run instr --updates 3000
python scripts/twostop_clock.py --clock 600 --per-spec 8 --out clock_600.json
python scripts/ghost_fleet.py --run myrun --suite medium --task 1 --rollouts 16
python scripts/paper_figures.py
\end{verbatim}
\end{quote}

The instrument channels are selected by the \code{SPACENAV\_INSTRUMENTS}
environment variable, which accepts a named subset, a four-bit mask, or
\code{0}/\code{1}. A checkpoint must be evaluated with the same setting it was
trained with; the observation shape is identical either way, so a mismatch
produces a policy reading zeros where it expects an instrument, with no error.

\section{Limitations and outlook}
\label{sec:limits}

The most substantive limitation is that we do not know \emph{why} redundant
observation features cost a policy ten points. Two mechanisms --- capacity
pressure and the conditioning of heavy-tailed ratio features --- predict
different behaviour under scale, and distinguishing them requires training at
embedding widths we have not been able to afford at five seeds per arm.

Second, the environment is not chaotic, and we would like it to be. Section
\ref{sec:chaos} suggests the route: closer encounters, tighter binaries, and
above all episodes that are long relative to orbital periods. That requires a
per-mission episode clock, which the current implementation does not have.

Third, our difficulty labels are only as good as the reference optimiser.
A mission labelled \emph{opaque} is one nothing we have tried can fly, which is a
statement about our search as much as about the mission.

What we think transfers beyond this particular environment is the shape of the
problem. A task can be cheap to simulate, deterministic, fully specified, and
still resist a one-dimensional notion of difficulty; the measurement of that
difficulty can be fragile enough that a published conclusion inverts on a larger
sample; and the information given to a learner about its own task can hurt it in
ways that depend not on whether the information is \emph{true} but on whether it
is \emph{new}. None of those three is specific to gravity.

\section*{Acknowledgements}

Development of this environment, including the experimental design, measurement
code and analysis reported here, was carried out in collaboration with Claude
(Anthropic), operating as an autonomous agent with access to the repository and
compute. The decision records referenced in Section~\ref{sec:repro} were written
during that process and include the agent's own retracted conclusions.

\begin{thebibliography}{99}
\small

\bibitem[Agarwal et al.(2021)]{agarwal2021}
R.~Agarwal, M.~Schwarzer, P.~S. Castro, A.~Courville, and M.~G. Bellemare.
\newblock Deep reinforcement learning at the edge of the statistical precipice.
\newblock \emph{Advances in Neural Information Processing Systems}, 2021.
\newblock arXiv:2108.13264.

\bibitem[Chirikov(1979)]{chirikov1979}
B.~V. Chirikov.
\newblock A universal instability of many-dimensional oscillator systems.
\newblock \emph{Physics Reports}, 52(5):263--379, 1979.

\bibitem[Colas et al.(2018)]{colas2018}
C.~Colas, O.~Sigaud, and P.-Y. Oudeyer.
\newblock How many random seeds? Statistical power analysis in deep
reinforcement learning experiments.
\newblock arXiv:1806.08295, 2018.

\bibitem[Daza et al.(2016)]{daza2016}
A.~Daza, A.~Wagemakers, B.~Georgeot, D.~Gu\'ery-Odelin, and M.~A.~F. Sanju\'an.
\newblock Basin entropy: a new tool to analyze uncertainty in dynamical systems.
\newblock \emph{Scientific Reports}, 6:31416, 2016.

\bibitem[Deck et al.(2013)]{deck2013}
K.~M. Deck, M.~Payne, and M.~J. Holman.
\newblock First-order resonance overlap and the stability of close two-planet
systems.
\newblock \emph{The Astrophysical Journal}, 774(2):129, 2013. arXiv:1307.8119.

\bibitem[Franklin et al.(2019)]{franklin2019}
D.~W. Franklin, S.~Franklin, and D.~M. Wolpert.
\newblock Fractionation of the visuomotor feedback response to directions of
movement and perturbation.
\newblock \emph{Journal of Neuroscience}, 2019.

\bibitem[Furuta et al.(2021)]{furuta2021}
H.~Furuta, T.~Matsushima, T.~Kozuno, Y.~Matsuo, S.~Levine, O.~Nachum, and
S.~S. Gu.
\newblock Policy information capacity: information-theoretic measure for task
complexity in deep reinforcement learning.
\newblock \emph{International Conference on Machine Learning}, 2021.

\bibitem[Grebogi et al.(1983)]{grebogi1983}
C.~Grebogi, S.~W. McDonald, E.~Ott, and J.~A. Yorke.
\newblock Final state sensitivity: an obstruction to predictability.
\newblock \emph{Physics Letters A}, 99(9):415--418, 1983.

\bibitem[Izzo and \"Ozt\"urk(2021)]{izzo2021}
D.~Izzo and E.~\"Ozt\"urk.
\newblock Real-time guidance for low-thrust transfers using deep neural
networks.
\newblock \emph{Journal of Guidance, Control, and Dynamics}, 44(2), 2021.

\bibitem[Lai and Gershman(2024)]{lai2024}
L.~Lai and S.~J. Gershman.
\newblock Human decision making balances reward maximization and policy
compression.
\newblock \emph{PLOS Computational Biology}, 2024.

\bibitem[Matthews et al.(2024)]{kinetix2024}
M.~Matthews, M.~Beukman, B.~Ellis, M.~Samvelyan, M.~Jackson, S.~Coward, and
J.~Foerster.
\newblock Kinetix: investigating the training of general agents through
open-ended physics-based control tasks.
\newblock arXiv:2410.23208, 2024.

\bibitem[Lee et al.(2025)]{simbav2}
H.~Lee, D.~Hwang, D.~Kim, H.~Kim, J.~J. Tai, K.~Subramanian, P.~R. Wurman,
J.~Choo, P.~Stone, and T.~Seno.
\newblock SimbaV2: hyperspherical normalization for scalable deep reinforcement
learning.
\newblock arXiv:2502.15280, 2025.

\bibitem[Repp(2005)]{repp2005}
B.~H. Repp.
\newblock Sensorimotor synchronization: a review of the tapping literature.
\newblock \emph{Psychonomic Bulletin \& Review}, 12(6):969--992, 2005.

\bibitem[Schulman et al.(2017)]{schulman2017}
J.~Schulman, F.~Wolski, P.~Dhariwal, A.~Radford, and O.~Klimov.
\newblock Proximal policy optimization algorithms.
\newblock arXiv:1707.06347, 2017.

\bibitem[Wang et al.(2023)]{wang2023}
A.~Wang, T.~Herbert, and S.~Gao.
\newblock On the fractal nature of the reinforcement learning objective
landscape.
\newblock \emph{Advances in Neural Information Processing Systems}, 2023.

\bibitem[Wisdom(1980)]{wisdom1980}
J.~Wisdom.
\newblock The resonance overlap criterion and the onset of stochastic behavior
in the restricted three-body problem.
\newblock \emph{The Astronomical Journal}, 85:1122--1133, 1980.

\end{thebibliography}

\end{document}
